% Part: turing-machines
% Chapter: machines-computations
% Section: variants

\documentclass[../../../include/open-logic-section]{subfiles}

\begin{document}

\olfileid{tur}{mac}{var}
\olsection{টুরিং যন্ত্রের বিভিন্ন রূপ}

টুরিং যন্ত্র সংজ্ঞায়িত করার বহু সম্ভাব্য উপায় আছে; আমাদেরটি তার মাত্র
একটি। কিছু দিক থেকে আমাদের সংজ্ঞা অন্যগুলির চেয়ে বেশি উদার। আমরা যেকোনো
সসীম বর্ণমালা অনুমোদন করি; আরও সীমিত কোনো সংজ্ঞায় কেবল দুটি ফিতা-প্রতীক,
$\TMstroke$ ও~$\TMblank$, অনুমোদিত হতে পারে। আমরা যন্ত্রকে একই সঙ্গে
ফিতায় প্রতীক লেখা ও চলার অনুমতি দিই; অন্য সংজ্ঞাগুলি লেখা অথবা চলা—দুটির
একটি অনুমোদন করে। আমরা ফিতার হেড না সরিয়েও লেখার সম্ভাবনা রাখি; অন্য
সংজ্ঞাগুলি $\TMstay$ ``নির্দেশটি'' বাদ দেয়। আবার কিছু দিক থেকে আমাদের
সংজ্ঞা বেশি সীমাবদ্ধ। আমরা ধরে নিয়েছি ফিতাটি কেবল এক দিকে অসীম; অন্য
সংজ্ঞায় ফিতাটি বাঁ ও ডান উভয় দিকেই অসীম। এমনকি যেকোনো সংখ্যক আলাদা
ফিতা, অথবা ঘরের একটি অসীম জালিকাও অনুমোদন করা যায়। আমরা টুরিং যন্ত্রের
নির্দেশ-সেটকে অবস্থান্তর অপেক্ষক দিয়ে প্রকাশ করি; অন্য সংজ্ঞায় অবস্থান্তর
সম্বন্ধ ব্যবহৃত হয়, যেখানে কোনো প্রদত্ত পরিস্থিতিতে যন্ত্রটির একাধিক
সম্ভাব্য নির্দেশ থাকে।

সংজ্ঞার শেষের এই শিথিলতাটি বিশেষভাবে আকর্ষণীয়। আমাদের সংজ্ঞায় যন্ত্রটি
দশা~$q$-তে থেকে প্রতীক~$\sigma$ পড়লে $\delta(q, \sigma)$ নতুন প্রতীক,
দশা ও ফিতা-হেডের অবস্থান নির্ধারণ করে। কিন্তু নির্দেশ-সেটকে যদি বর্তমান
দশা-প্রতীক জোড়া $\tuple{q, \sigma}$ এবং নতুন দশা-প্রতীক-দিক ত্রয়ী
$\tuple{q', \sigma', D}$-র মধ্যেকার সম্বন্ধ হিসেবে ধরি, তবে টুরিং
যন্ত্রের কাজটি এককভাবে নির্ধারিত নাও হতে পারে—নির্দেশ-সম্বন্ধে
$\tuple{q, \sigma, q', \sigma', D}$ ও
$\tuple{q, \sigma, q'', \sigma'', D'}$ দুটিই থাকতে পারে। সেক্ষেত্রে
আমরা একটি \emph{অনির্ধারণবাদী} টুরিং যন্ত্র পাই। গণনামূলক জটিলতা তত্ত্বে
এগুলির গুরুত্বপূর্ণ ভূমিকা আছে।

টুরিং যন্ত্র কখন থামে, তা নিয়েও বিভিন্ন রীতি আছে: আমাদের মতে অবস্থান্তর
অপেক্ষক অসংজ্ঞায়িত হলে সেটি থামে; অন্য সংজ্ঞায় যন্ত্রটিকে বিশেষ কোনো
নির্দিষ্ট থামা-দশায় থাকতে হয়। নির্দিষ্ট থামা-দশা আবশ্যক করলে টুরিং
যন্ত্র কী গণনা করতে পারে তার উপর কেন কোনো সীমা পড়ে না, তা আমরা
\olref[hal]{sec}-এ ব্যাখ্যা করেছি। আমাদের টুরিং যন্ত্রের ফিতা কেবল এক
দিকে অসীম বলে কোনো কোনো ক্ষেত্রে যন্ত্রটি একটি নির্দেশ যথাযথভাবে পালন
করতে পারে না: যেমন, বাঁদিকের প্রথম ঘরটি পড়ার সময় যদি তাকে বাঁ দিকে
যেতে বলা হয়। আমাদের সংজ্ঞা অনুযায়ী ``বাইরে পড়ে যাওয়ার'' বদলে সেটি
একই জায়গায় থাকে; কিন্তু এমন ঘটলে যন্ত্রটি থামবে—এভাবেও সংজ্ঞা দেওয়া
যেত। এই সংজ্ঞাটিও সমতুল্য। ঘর~$0$-এর সীমা-চিহ্নটি অক্ষুণ্ণ থাকলে,
ঘর~$0$ থেকে বাঁ দিকে যাওয়ার চেষ্টা করলেই যে টুরিং যন্ত্র থামে তার আচরণ
আমরা প্রতিটি অবস্থান্তর $\delta(q,\TMendtape) =
\tuple{q',\sigma,\TMleft}$ মুছে দিয়ে অনুকরণ করতে পারি—তাহলে
$\TMendtape$-এ বাঁ দিকে যাওয়ার চেষ্টা না করে যন্ত্রটি থেমে যায়।
\footnote{এটি \emph{পুরোপুরি} কাজ করে না, কারণ ঘর~$0$ ছাড়া অন্য ঘরেও
$\TMendtape$ লেখা ও পড়া থেকে আমাদের কিছুই আটকায় না
(\olref[una]{ex:mover} দেখো), আবার ঘর~$0$-এর চিহ্নটিও মুছে ফেলা যায়।
ঘর~$0$-এর জন্য দ্বিতীয় একটি~$\TMendtape'$ প্রতীক সংরক্ষণ করে, সেটিকে
কখনও না মুছে এবং অন্য কোথাও না লিখে সমস্যাটি এড়ানো যায়; বাঁমুখী
অবস্থান্তরগুলি তখন ওই সংরক্ষিত চিহ্নের উপরেই মুছতে হবে।}
% BN-SRC-174 TARGET-CORRECTION-BEGIN
\par\smallskip\noindent\textnormal{\small\textbf{উৎস-সংশোধন (BN-SRC-174):} হিমায়িত ইংরেজি নির্মাণটি শুধু অন্য ঘরে শেষ-চিহ্ন ব্যবহারের অসুবিধা স্বীকার করে। কিন্তু তার নিজের উদার সংজ্ঞায় ঘর~$0$-এর চিহ্নও মোছা যায়; তখন শুধু শেষ-চিহ্ন পড়ে বাঁ দিকে যাওয়া অবস্থান্তর মুছলে ভৌত বাঁ সীমানা শনাক্ত হয় না। বাংলা পাঠে নতুন চিহ্নটিকে শুধু ঘর~$0$-এর জন্য সংরক্ষিত ও অমোচনীয় করার পূর্ণ শর্তটি দেওয়া হয়েছে।} % BN-SRC-174 TARGET-CORRECTION-END

সংখ্যা প্রকাশের (এবং তাই টুরিং যন্ত্রে গণিত নিবেশ-নির্গম অপেক্ষকটি
প্রকাশের) বিভিন্ন উপায়ও আছে: আমরা এককীয় উপস্থাপনা ব্যবহার করি, কিন্তু
দ্বিমিক উপস্থাপনাও ব্যবহার করা যায়। তার জন্য $\TMblank$ ও~$\TMendtape$-
এর অতিরিক্ত দুটি প্রতীক লাগে।

এবার একটি আকর্ষণীয় তথ্য: কোন অপেক্ষকগুলি টুরিং-গণনসাধ্য, তার বিচারে এই
রূপভেদগুলির কোনোটি গুরুত্বপূর্ণ নয়। \emph{কোনো একটি সংজ্ঞা অনুসারে একটি
অপেক্ষক টুরিং-গণনসাধ্য হলে সব সংজ্ঞা অনুসারেই সেটি টুরিং-গণনসাধ্য।}

এটি যাচাই করার খুঁটিনাটিতে আমরা যাব না। শুধু একটি উদাহরণ দেখা যাক: উভয়
দিকে অসীম ফিতা বা একাধিক ফিতা অনুমোদন করলেও অতিরিক্ত গণনক্ষমতা পাওয়া
যায় না। মোটামুটি কারণটি হলো, একটিমাত্র একমুখী অসীম ফিতাযুক্ত টুরিং যন্ত্র
একাধিক অথবা দ্বিমুখী অসীম ফিতা অনুকরণ করতে পারে। যেমন, অতিরিক্ত দশা ও
নির্দেশ ব্যবহার করে একাধিক ফিতা বা দ্বিমুখী অসীম ফিতাযুক্ত যন্ত্রের একটি
প্রোগ্রামকে একটিমাত্র একমুখী অসীম ফিতাযুক্ত যন্ত্রের প্রোগ্রামে
``অনুবাদ'' করা যায়। অনূদিত যন্ত্রটি জোড় ঘরগুলি ফিতা~$1$-এর ঘর হিসেবে
(অথবা দ্বিমুখী অসীম ফিতার ``ধনাত্মক'' ঘর হিসেবে) এবং বিজোড় ঘরগুলি
ফিতা~$2$-এর ঘর হিসেবে (অথবা ``ঋণাত্মক'' ঘর হিসেবে) ব্যবহার করতে পারে।

\end{document}
